% !TeX program = xelatex
% !TeX encoding = UTF-8



\documentclass[UTF8,a4paper,zihao=-4,fontset=fandol]{ctexart}

\usepackage{geometry}
\usepackage{graphicx}
\usepackage{booktabs}
\usepackage{tabularx}
\usepackage{array}
\usepackage{amsmath}
\usepackage{tikz}
\usepackage{caption}
\usepackage[hidelinks]{hyperref}
\usepackage{enumitem}
\usepackage{float}
\usepackage{fancyhdr}
\usepackage{setspace}

\usetikzlibrary{arrows.meta,positioning,fit}

% 页面设置
\geometry{left=2.8cm,right=2.8cm,top=2.6cm,bottom=2.6cm}

% 正文段落设置
\setlength{\parindent}{2em}
\setlength{\parskip}{0pt}
\setstretch{1.3}

% 目录只显示到二级标题，避免目录过长、过散
\setcounter{tocdepth}{2}

% 图表与正文间距，避免表前表后、图前图后过松
\setlength{\intextsep}{8pt plus 1pt minus 1pt}
\setlength{\textfloatsep}{8pt plus 1pt minus 1pt}
\setlength{\floatsep}{8pt plus 1pt minus 1pt}
\setlength{\abovecaptionskip}{4pt}
\setlength{\belowcaptionskip}{2pt}

% 列表间距
\setlist{
  itemsep=0pt,
  topsep=2pt,
  parsep=0pt,
  partopsep=0pt
}

% 字体设置
% Windows 本地编译时优先使用宋体、黑体、楷体；
% Overleaf 或无对应字体环境下使用 Fandol 字体，保证可编译。
\IfFontExistsTF{SimSun}{\setCJKmainfont{SimSun}}{}
\IfFontExistsTF{SimHei}{\setCJKsansfont{SimHei}}{}
\IfFontExistsTF{KaiTi}
  {\setCJKfamilyfont{kai}{KaiTi}}
  {\setCJKfamilyfont{kai}{FandolKai-Regular}}
\IfFontExistsTF{FangSong}
  {\setCJKfamilyfont{fs}{FangSong}}
  {\setCJKfamilyfont{fs}{FandolFang-Regular}}
\IfFontExistsTF{FZXiaoBiaoSong-Z05S}
  {\setCJKfamilyfont{xbs}{FZXiaoBiaoSong-Z05S}}
  {\setCJKfamilyfont{xbs}{FandolSong-Regular}}

\newcommand{\xbs}{\CJKfamily{xbs}}

\IfFontExistsTF{Times New Roman}{\setmainfont{Times New Roman}}{}

% 标题格式：标题黑体，正文宋体
\ctexset{
  section = {
    format = \heiti\zihao{-3}\bfseries,
    name = {,、},
    number = \chinese{section},
    aftername = \hspace{0pt},
    beforeskip = 1.0em,
    afterskip = 0.6em
  },
  subsection = {
    format = \heiti\zihao{4}\bfseries,
    beforeskip = 0.8em,
    afterskip = 0.4em
  },
  subsubsection = {
    format = \heiti\zihao{-4}\bfseries,
    beforeskip = 0.6em,
    afterskip = 0.3em
  },
  paragraph = {
    format = \heiti\zihao{-4},
    beforeskip = 0.5em,
    afterskip = 0.2em
  }
}

% 图表标题格式
\DeclareCaptionFont{songcap}{\songti\zihao{5}}
\captionsetup[figure]{
  font=songcap,
  justification=centering,
  labelsep=quad,
  skip=4pt
}
\captionsetup[table]{
  font=songcap,
  justification=centering,
  labelsep=quad,
  skip=4pt
}

% 页码
\pagestyle{fancy}
\fancyhf{}
\fancyfoot[C]{\songti\zihao{-4}\thepage}
\renewcommand{\headrulewidth}{0pt}
\renewcommand{\footrulewidth}{0pt}
\fancypagestyle{plain}{
  \fancyhf{}
  \fancyfoot[C]{\songti\zihao{-4}\thepage}
  \renewcommand{\headrulewidth}{0pt}
  \renewcommand{\footrulewidth}{0pt}
}

\renewcommand{\contentsname}{\heiti\zihao{4} 目录}
\renewcommand{\refname}{\heiti\zihao{4} 参考文献}

% 表格行距，不要太松
\renewcommand{\arraystretch}{1.08}

\hypersetup{
  colorlinks=false,
  pdfborder={0 0 0}
}

% 图片命令：统一固定位置，避免图片跑到标题前面
\newcommand{\placeholderfigure}[3]{
\begin{figure}[H]
  \centering
  \IfFileExists{#1}{
    \includegraphics[width=0.86\textwidth]{#1}
  }{
    \fbox{
      \begin{minipage}[c][0.38\textwidth][c]{0.86\textwidth}
        \centering
        \songti\zihao{-4}
        图片占位：\texttt{#1}\\
        请将对应页面截图放入 \texttt{figures/} 目录后重新编译。
      \end{minipage}
    }
  }
  \caption{#2}
  \label{#3}
\end{figure}
}

\title{\xbs\zihao{3}《安养智巡》\\[0.45em]
{\heiti\zihao{4} 康养机构AI安全风险预警系统}\\[0.25em]
{\songti\zihao{-4} 项目研究报告}}
\author{}
\date{}

\begin{document}
\songti\zihao{-4}
\setstretch{1.3}

\maketitle

\phantomsection
\addcontentsline{toc}{section}{摘要}
\begin{center}
{\heiti\zihao{4} 摘要}
\end{center}

安养智巡是面向老年人疗养中心的后台管理系统，围绕健康风险、心理风险和安全风险的早发现、快研判、可追踪处置构建多模态 AI 融合预警体系。系统以摄像头视频、床垫传感器、门磁、毫米波雷达、可穿戴设备、护理记录、老人健康档案、行为画像、历史事件记录和应急知识库为数据基础，结合 YOLO 目标检测、YOLO 姿态估计、行为偏离度分析、风险评分、规则引擎和知识库辅助决策，实现风险识别、智能研判、派单处置、应急流程、事件归档和数据看板的闭环管理。

系统已完成原型部署与模拟运行验证，后端服务接口、视觉识别服务、规则研判服务、数据看板和业务闭环模块均已形成可运行能力。在测试和模拟运行中，系统可呈现今日告警 26 起、待处理高风险 7 起、平均响应 4.2 分钟、设备在线率 97.6\%、摄像头在线 36/38、传感器在线 124/128、监测延迟 142 ms 等运行指标。项目重点不是替代护理员和医生，而是通过连续感知、跨模态校验和可解释研判，辅助机构提升异常发现效率、处置一致性和护理管理精细化水平。

\noindent{\heiti\zihao{-4} 关键词：}智慧养老；多模态融合预警；YOLO；姿态估计；行为画像；护理决策支持；可解释人工智能

\newpage
\begingroup
\setstretch{1.05}
\tableofcontents
\endgroup
\newpage

\section{项目的背景和国内外研究现状}

\subsection{项目背景}

我国人口老龄化程度持续加深，疗养中心、护理院和医养结合机构承担着越来越复杂的照护任务。老年人群体常见风险并不只表现为跌倒、烟火异常、无人看护等显性安全事件，还包括久卧未动、夜间离床频繁、睡眠异常、活动量下降、情绪低落、孤独风险、焦虑风险、慢性病指标波动等健康与心理风险。这些风险往往具有隐匿性、累积性和个体差异，如果仅依赖人工巡房、纸质记录和事后上报，容易出现发现滞后、处置链条断裂和护理经验难以沉淀等问题。

智慧养老研究表明，传感器、物联网、人工智能和数字平台可以提升养老服务的连续监测能力与管理效率，但真实机构场景中仍存在数据来源分散、预警误报率高、护理处置与系统告警脱节等问题。因此，面向疗养中心后台管理的系统需要同时满足三类要求：第一，能够接入多源数据并进行稳定融合；第二，能够对健康、心理和安全风险进行可解释研判；第三，能够与护理排班、派单、应急流程、事件归档和管理看板形成闭环。

\subsection{国内外研究现状}

\subsubsection{智慧养老与机构数字化管理研究}

关于中国智慧居家养老的发展研究指出，政策推动、平台建设和服务整合促进了智慧养老快速发展，但技术适配、需求匹配和持续运营仍是落地过程中的关键问题\cite{zhang2020smart}。从智慧养老服务体系角度看，智能感知、平台化服务和健康管理已经成为养老服务升级的重要方向\cite{hung2023smart}。关于 IoT 可穿戴设备在老年照护中的综述研究进一步表明，可穿戴设备、智能家居、环境传感器和健康监测设备已经成为老年照护技术的重要组成部分，但互操作性、隐私保护和长期使用接受度仍是关键挑战\cite{stavropoulos2020iot}。另有研究指出，可穿戴技术可用于室内定位、活动识别和生命体征实时监测，为机构照护中的连续感知提供了重要技术支撑\cite{dong2017wearable}。

现有研究为本项目提供了重要基础：养老机构需要的不只是单点设备，而是能够把传感、记录、研判和处置组织起来的后台管理系统。安养智巡因此采用机构级后台架构，将多模态感知、护理业务流程和管理看板统一设计。

\subsubsection{跌倒检测、人体姿态估计与行为识别研究}

跌倒检测是老年照护场景中最受关注的安全风险之一。跌倒检测综述研究将相关系统归纳为基于可穿戴传感器、环境传感器和视觉感知等多种路线\cite{mubashir2013survey}。另有研究指出，跌倒检测系统的难点在于真实跌倒样本稀缺、日常动作与跌倒动作相似、不同环境下模型泛化不足等\cite{igual2013challenges}。这说明单一视觉或单一传感器判断很难覆盖疗养中心的复杂场景。

在视觉识别方面，YOLO 将目标检测转化为端到端回归问题，显著提升了实时检测效率\cite{redmon2016yolo}。OpenPose 通过 Part Affinity Fields 实现多人二维姿态估计，为人体姿态理解、跌倒识别和异常行为分析提供了重要技术基础\cite{cao2017realtime}。Ultralytics YOLO11 文档显示，YOLO11 支持检测、分割、分类、姿态估计等任务，并包含轻量化模型配置\cite{ultralytics2024yolo11}。本项目在系统设计中接入 YOLO11n 和 YOLO11n-pose，用于人员、床位、异常姿态、疑似跌倒、长时间静止等场景识别。

\subsubsection{多模态融合、健康监测与心理风险评估研究}

多模态机器学习关注不同数据模态之间的表示、对齐、融合和协同学习。多模态机器学习分类框架指出，融合不仅是简单拼接，而涉及表示学习、模态对齐、跨模态补全和融合决策等核心问题\cite{baltrusaitis2019multimodal}。关于多模态机器学习基础问题的研究进一步强调，异构数据、时序对齐、缺失模态和可靠性评估是多模态系统落地的重要难点\cite{liang2022foundations}。

在健康和心理风险方面，孤独与社会隔离相关研究表明，孤独和社会隔离与死亡风险升高相关，提示养老服务不能只关注生理安全，也需要关注社会连接和心理状态\cite{holtlunstad2015loneliness}。数字健康工具用于抑郁被动监测的系统综述说明，基于手机、可穿戴设备和行为数据的被动监测可为心理健康风险评估提供新路径\cite{sequeira2021digital}。这些研究启发本项目将夜间离床频次、活动量下降、久卧未动、互动减少、睡眠异常和护理记录纳入行为画像，用于识别健康和心理风险趋势。

\subsubsection{护理决策支持与可解释人工智能研究}

护理场景中的智能系统不能只输出告警，还需要给出可理解、可追溯、可执行的处置建议。临床决策支持系统相关研究表明，面向个体的建议、与工作流集成和自动触发是提升系统有效性的关键因素\cite{kawamoto2005cdss}。相关系统综述也指出，临床决策支持系统通常能改善医疗服务过程，但效果取决于系统设计、嵌入时机和实际使用方式\cite{bright2012cdss}。可解释人工智能研究强调，复杂模型需要向用户提供可理解的理由、证据和不确定性信息，才能支持人机协同决策\cite{gunning2019xai}。SHAP 方法则为特征贡献度解释提供了统一思路\cite{lundberg2017shap}。

综上，国内外研究已经在智慧养老、传感器监测、跌倒检测、姿态估计、心理健康监测、护理决策支持和可解释 AI 等方面形成了基础，但仍存在以下不足：第一，单一模态系统容易误报或漏报；第二，许多系统过度关注跌倒等显性事件，对健康和心理风险关注不足；第三，预警与护理处置、事件归档、统计复盘之间存在脱节；第四，风险研判可解释性不足；第五，缺少适合疗养中心后台管理的机构级闭环系统。安养智巡正是在这些不足基础上，提出面向疗养中心的多模态融合预警与护理业务闭环方案。

\section{项目研究的内容和技术路线}

\subsection{总体研究内容}

安养智巡围绕“感知、研判、处置、复盘”四个环节展开研究和实现。系统通过摄像头视频/图像数据、床垫传感器、门磁、毫米波雷达、可穿戴设备、护理记录、老人健康档案、行为画像、历史事件记录和应急知识库构建多源数据底座；通过 YOLO 目标检测、YOLO 姿态估计、区域目标识别和异常行为识别形成实时感知结果；通过行为偏离度分析、多模态风险评分和规则引擎生成风险等级、触发原因、置信度和建议动作；最后通过派单调度、应急流程、事件管理和数据看板形成护理管理闭环。

\begin{figure}[H]
  \centering
  \begin{tikzpicture}[
    node distance=0.8cm and 0.7cm,
    box/.style={draw, rounded corners=2pt, minimum width=3.4cm, minimum height=0.72cm, align=center, font=\songti\zihao{5}},
    core/.style={draw, rounded corners=2pt, minimum width=4.0cm, minimum height=0.86cm, align=center, font=\songti\zihao{5}},
    arrow/.style={-{Stealth[length=2.0mm]}, line width=0.5pt}
  ]
    \node[box] (device) {多源感知设备};
    \node[box, right=of device] (service) {后端服务接口};
    \node[core, right=of service] (fusion) {多模态融合预警};
    \node[box, below=of service] (knowledge) {规则引擎与护理知识库};
    \node[box, right=of fusion] (business) {派单处置与事件归档};
    \node[box, below=of business] (dashboard) {数据看板与管理复盘};
    \draw[arrow] (device) -- (service);
    \draw[arrow] (service) -- (fusion);
    \draw[arrow] (knowledge) -- (fusion);
    \draw[arrow] (fusion) -- (business);
    \draw[arrow] (business) -- (dashboard);
    \draw[arrow] (dashboard) -- (knowledge);
  \end{tikzpicture}
  \caption{系统总体架构图}
  \label{fig:system-architecture}
\end{figure}

如图 \ref{fig:system-architecture} 所示，系统以多源感知设备和后端接口为基础，将视觉识别、传感器状态、健康档案和护理记录统一汇入融合预警模块，再由规则引擎和护理知识库生成可解释处置建议，最终进入派单处置、事件归档和数据看板。

系统的数据来源和业务用途如表 \ref{tab:data-sources} 所示。表中将设备侧数据、护理业务数据和历史画像数据放在同一结构下说明，便于体现多模态融合并不是单一视觉识别，而是围绕健康、心理和安全风险共同建模。

\begin{table}[H]
  \centering
  \caption{多模态数据来源与业务用途}
  \label{tab:data-sources}
  \zihao{5}
  \begin{tabularx}{\textwidth}{>{\centering\arraybackslash}p{3.2cm}>{\centering\arraybackslash}p{4.2cm}>{\centering\arraybackslash}X}
    \toprule
    数据来源 & 典型数据 & 主要用途 \\
    \midrule
    摄像头视频/图像 & 人员位置、姿态、区域目标、异常场景 & 跌倒识别、长时间静止识别、无人看护和冲突场景识别 \\
    床垫传感器 & 在床状态、离床时间、卧床时长、睡眠片段 & 夜间离床、久卧未动、睡眠异常和回床确认 \\
    门磁与毫米波雷达 & 出入记录、房间活动、静止状态、区域存在 & 离房未归、区域滞留、弱光或遮挡场景补充判断 \\
    可穿戴设备 & 心率、活动量、步数、部分健康指标 & 健康风险趋势、活动量下降和异常体征辅助判断 \\
    护理记录与健康档案 & 慢病情况、护理等级、既往事件、医嘱摘要 & 个体化风险阈值、处置建议和医生复核依据 \\
    行为画像与历史事件 & 作息规律、异常频次、响应记录、处置结果 & 行为偏离度计算、风险趋势预测和管理复盘 \\
    \bottomrule
  \end{tabularx}
\end{table}

在数据接入和风险研判基础上，系统主要功能模块如表 \ref{tab:functions} 所示。各模块与后台页面一一对应，并共同服务于风险识别、处置调度、事件归档和管理复盘。

\begin{table}[H]
  \centering
  \caption{系统主要功能模块}
  \label{tab:functions}
  \zihao{5}
  \begin{tabularx}{\textwidth}{>{\centering\arraybackslash}p{3cm}>{\centering\arraybackslash}X>{\centering\arraybackslash}p{3.2cm}}
    \toprule
    模块 & 主要功能 & 对应页面 \\
    \midrule
    首页运行总览 & 汇总今日告警、高风险事件、平均响应、设备在线率和待办任务 & 首页运行总览 \\
    实时监测 & 接入摄像头画面、视觉识别结果、模型状态和实时事件队列 & 实时监测页面 \\
    健康监护 & 聚合身体指标、心理风险、近期事件、睡眠与护理建议 & 健康监护页面 \\
    风险调度 & 根据风险等级、等待时长、护理员状态和区域位置进行优先级排序与派单 & 风险调度页面 \\
    行为画像 & 分析作息、活动区域、夜间离床、久卧未动和行为偏离度 & 行为画像页面 \\
    应急流程 & 为跌倒、烟火、无人看护、心理风险等事件生成处置步骤和注意事项 & 应急流程页面 \\
    事件管理 & 完成事件查询、处置记录、归档、复盘和责任追踪 & 事件管理页面 \\
    数据看板 & 统计风险趋势、响应效率、设备在线率、区域分布和护理负载 & 数据看板页面 \\
    知识库工作台 & 管理护理知识、规则依据、问答内容和处置建议模板 & 知识库工作台页面 \\
    \bottomrule
  \end{tabularx}
\end{table}

\subsection{核心技术路线}

系统技术路线分为六层：数据采集层、感知识别层、行为分析层、多模态融合预警层、智能决策层和业务闭环层。其逻辑关系如图 \ref{fig:tech-route} 所示。

\begin{figure}[H]
  \centering
  \begin{tikzpicture}[
    node distance=0.72cm,
    layer/.style={draw, rounded corners=2pt, minimum width=0.86\textwidth, minimum height=0.78cm, align=center, font=\songti\zihao{5}},
    arrow/.style={-{Stealth[length=2.2mm]}, line width=0.5pt}
  ]
    \node[layer] (data) {数据采集层：摄像头、床垫传感器、门磁、毫米波雷达、可穿戴设备、护理记录、健康档案};
    \node[layer, below=of data] (perception) {感知识别层：YOLO 目标检测、YOLO 姿态估计、人体姿态识别、区域目标识别、异常行为识别};
    \node[layer, below=of perception] (behavior) {行为分析层：历史作息、夜间离床、活动区域、近期异常、健康指标、护理记录与行为偏离度};
    \node[layer, below=of behavior] (fusion) {多模态融合预警层：风险评分、风险等级、风险原因、置信度、可解释因子};
    \node[layer, below=of fusion] (decision) {智能决策层：规则引擎、护理知识库、处置建议、注意事项、应急流程、引用依据};
    \node[layer, below=of decision] (closedloop) {业务闭环层：事件队列、优先级排序、派单、处置、记录、归档、统计分析};
    \draw[arrow] (data) -- (perception);
    \draw[arrow] (perception) -- (behavior);
    \draw[arrow] (behavior) -- (fusion);
    \draw[arrow] (fusion) -- (decision);
    \draw[arrow] (decision) -- (closedloop);
  \end{tikzpicture}
  \caption{安养智巡技术路线图}
  \label{fig:tech-route}
\end{figure}

\subsection{多模态风险评分模型}

系统将视觉识别、传感器状态、健康指标、行为画像、心理风险和历史事件统一映射为风险特征。设老人或场景对象为 $i$，视觉风险特征为 $V_i$，传感器风险特征为 $S_i$，健康指标风险特征为 $H_i$，行为偏离特征为 $B_i$，历史事件特征为 $E_i$，则综合风险评分可表示为：
\begin{equation}
R_i=\sigma(w_vV_i+w_sS_i+w_hH_i+w_bB_i+w_eE_i+b)
\label{eq:risk-score}
\end{equation}
其中，$w_v,w_s,w_h,w_b,w_e$ 为各类特征权重，$b$ 为偏置项，$\sigma(\cdot)$ 为归一化函数。系统根据 $R_i$ 将风险划分为低风险、中风险、高风险和紧急风险，并输出触发因子与处置建议。

在调度环节，系统还综合风险等级、等待时长、护理员可用状态和空间距离计算任务优先级：
\begin{equation}
P_i=\alpha R_i+\beta T_i+\gamma U_i+\delta L_i
\label{eq:priority}
\end{equation}
其中，$T_i$ 表示等待时长，$U_i$ 表示护理员负载与可用状态，$L_i$ 表示区域距离或到达成本，$\alpha,\beta,\gamma,\delta$ 为可配置权重。该模型使紧急事件、长时间未响应事件和近距离可处置事件能够被优先派发。

风险等级划分与处置策略如表 \ref{tab:risk-levels} 所示。系统将评分结果转化为护理人员能够理解和执行的等级规则，避免算法输出停留在抽象分值层面。

\begin{table}[H]
  \centering
  \caption{风险等级划分与处置策略}
  \label{tab:risk-levels}
  \zihao{5}
  \begin{tabularx}{\textwidth}{>{\centering\arraybackslash}p{2.4cm}>{\centering\arraybackslash}p{3.2cm}>{\centering\arraybackslash}X>{\centering\arraybackslash}p{3.2cm}}
    \toprule
    风险等级 & 典型触发条件 & 系统动作 & 处置要求 \\
    \midrule
    低风险 & 单项指标轻微偏离、短时活动减少 & 记录观察，进入健康趋势统计 & 由护理员在日常巡查中确认 \\
    中风险 & 多项信号持续异常、夜间离床频次升高 & 生成提醒，进入待办队列 & 护理员复核并补充护理记录 \\
    高风险 & 久卧未动、离床未归、健康指标异常叠加 & 提升优先级，触发派单与主管提醒 & 值班人员限时响应并记录结果 \\
    紧急风险 & 疑似跌倒未响应、烟火异常、严重冲突或无人看护 & 立即告警，启动应急流程 & 现场处置、上报主管并完成归档 \\
    \bottomrule
  \end{tabularx}
\end{table}

\subsection{系统运行数据与测试场景}

系统测试与模拟运行指标如表 \ref{tab:runtime-data} 所示。该组数据用于说明系统已能够围绕告警数量、高风险事件、响应效率和设备在线状态形成运行监测能力。

\begin{table}[H]
  \centering
  \caption{系统测试与模拟运行指标}
  \label{tab:runtime-data}
  \zihao{5}
  \begin{tabular}{ccc}
    \toprule
    指标 & 数值 & 业务含义 \\
    \midrule
    今日告警 & 26 起 & 当日由系统识别并进入队列的风险事件总量 \\
    待处理高风险 & 7 起 & 需要值班人员重点处理的风险事件 \\
    紧急事件 & 2 起 & 需要立即响应的高优先级事件 \\
    平均响应 & 4.2 分钟 & 从告警生成到人员响应的平均用时 \\
    设备在线率 & 97.6\% & 摄像头与传感器整体在线状态 \\
    摄像头在线 & 36/38 & 视频监测设备可用情况 \\
    传感器在线 & 124/128 & 床垫、门磁、毫米波雷达等设备可用情况 \\
    监测延迟 & 142 ms & 感知结果进入系统队列的平均延迟 \\
    \bottomrule
  \end{tabular}
\end{table}

系统覆盖的示例风险事件包括跌倒未响应、夜间离床未归、久卧未动、烟火异常、长时间滞留、无人看护、老人冲突、健康风险和心理风险。系统在设计上特别强调健康与心理风险，因为这类风险往往不是瞬时事件，而是通过多天行为变化、护理记录和身体指标逐渐显现。

典型风险事件与闭环处理方式如表 \ref{tab:event-closed-loop} 所示。表中列出的识别依据、处置建议和归档数据对应系统从预警到复盘的主要业务字段。

\begin{table}[H]
  \centering
  \caption{典型风险事件与闭环处理方式}
  \label{tab:event-closed-loop}
  \zihao{5}
  \begin{tabularx}{\textwidth}{>{\centering\arraybackslash}p{3.0cm}>{\centering\arraybackslash}p{3.5cm}>{\centering\arraybackslash}X>{\centering\arraybackslash}p{3.1cm}}
    \toprule
    事件类型 & 主要识别依据 & 处置建议 & 归档数据 \\
    \midrule
    跌倒未响应 & 姿态异常、地面区域停留、长时间无动作 & 立即派单，现场确认意识状态与外伤情况 & 图像证据、响应时长、处置结果 \\
    夜间离床未归 & 床垫空床、门磁出入、超过回床阈值 & 联系就近护理员巡查房间与走廊 & 离床时间、回床时间、巡查记录 \\
    久卧未动 & 长时间在床、活动量下降、护理记录异常 & 提醒翻身、压疮风险观察和健康复核 & 卧床时长、护理措施、后续观察 \\
    心理风险 & 互动减少、睡眠异常、活动量下降、历史记录叠加 & 提醒护理员访谈，必要时转医生复核 & 风险原因、访谈记录、医生建议 \\
    烟火异常 & 视觉烟火识别、区域告警、环境传感器异常 & 启动应急流程，通知主管并疏散相关区域 & 告警时间、处置人员、复盘结论 \\
    \bottomrule
  \end{tabularx}
\end{table}

\subsection{页面模块与实现效果}

登录页用于护理员、值班主管、机构管理者和医生等角色进入后台系统。页面根据账号角色控制功能入口和数据权限，保证不同岗位只访问其职责范围内的老人档案、事件记录、风险处置和管理统计数据，如图 \ref{fig:login} 所示。

\placeholderfigure{figures/login.png}{登录页效果图}{fig:login}

首页运行总览用于呈现机构当日运行状态，包括今日告警、待处理高风险、平均响应、设备在线率、摄像头在线数、传感器在线数和实时监测延迟。值班主管可以通过该页面快速判断当前机构风险压力、设备状态和护理响应效率，如图 \ref{fig:home-overview} 所示。

\placeholderfigure{figures/home-overview.png}{首页运行总览效果图}{fig:home-overview}

实时监测页面用于接入摄像头画面和视觉识别结果，展示目标检测框、姿态关键点、风险信号、模型版本、处理耗时和实时事件队列。该页面帮助值班人员快速定位疑似跌倒、长时间静止、区域滞留、无人看护和异常冲突等场景，并将高风险事件推送至调度流程，如图 \ref{fig:realtime-monitor} 所示。

\placeholderfigure{figures/realtime-monitor.png}{实时监测页面效果图}{fig:realtime-monitor}

健康监护页面以老人个体档案为中心，展示身体健康指标、心理健康风险、睡眠状态、近期事件和护理建议。页面将可穿戴设备、床垫传感器、护理记录和历史事件融合，用于识别久卧未动、活动量下降、夜间离床频繁、睡眠异常、情绪低落和孤独风险等非显性风险，如图 \ref{fig:health-monitor} 所示。

\placeholderfigure{figures/health-monitor.png}{健康监护页面效果图}{fig:health-monitor}

风险调度页面用于将事件队列转化为可执行护理任务。系统根据风险等级、等待时长、护理员状态、区域距离和事件类型进行优先级排序，支持派单、接单、处理中、完成和超时提醒等状态管理，提升高风险事件处置效率，如图 \ref{fig:dispatch} 所示。

\placeholderfigure{figures/dispatch.png}{风险调度页面效果图}{fig:dispatch}

行为画像页面用于分析老人历史作息、夜间离床次数、活动区域、近期异常次数、健康指标和护理记录。系统通过行为偏离度识别风险趋势，例如某位老人连续数日活动量下降并伴随夜间离床增多，系统会提示潜在健康或心理风险，而不是等到严重事件发生后才记录，如图 \ref{fig:behavior-profile} 所示。

\placeholderfigure{figures/behavior-profile.png}{行为画像页面效果图}{fig:behavior-profile}

应急流程页面针对跌倒未响应、烟火异常、夜间离床未归、无人看护、心理风险等事件生成处置步骤、注意事项和现场护理建议。页面将知识库依据与操作流程结合，帮助护理员在紧急情况下按步骤完成确认、呼叫、现场处置、上报和记录，如图 \ref{fig:emergency-flow} 所示。

\placeholderfigure{figures/emergency-flow.png}{应急流程页面效果图}{fig:emergency-flow}

事件管理页面用于完成风险事件的查询、筛选、处置记录、图片或识别结果追踪、责任人记录、归档和复盘。机构管理者可以通过事件管理页面追踪高发事件类型、重复发生区域和处置质量，为后续优化护理排班和设备布设提供依据，如图 \ref{fig:event-management} 所示。

\placeholderfigure{figures/event-management.png}{事件管理页面效果图}{fig:event-management}

数据看板页面用于统计风险趋势、事件类型分布、响应时长、设备在线率、区域风险热度和护理负载。该页面面向机构管理者，支持从日、周、月等维度观察机构运行状态，辅助管理者发现高风险区域和流程瓶颈，如图 \ref{fig:dashboard} 所示。

\placeholderfigure{figures/dashboard.png}{数据看板页面效果图}{fig:dashboard}

知识库工作台用于管理跌倒、久卧、夜间离床、心理风险、无人看护等事件的处置知识、注意事项和引用依据。系统在生成风险研判时可调用知识库内容，向护理员和医生输出可解释的处置建议，如图 \ref{fig:knowledge-workbench} 所示。

\placeholderfigure{figures/knowledge-workbench.png}{知识库工作台页面效果图}{fig:knowledge-workbench}

\section{项目的创新点}

\subsection{多模态融合预警机制}

项目将视觉识别、床垫传感器、门磁、毫米波雷达、可穿戴设备、护理记录、健康档案和行为画像统一纳入风险研判流程。与单一传感器预警相比，多模态融合能够通过跨模态校验降低误报。例如，视觉模型识别到老人疑似跌倒时，系统会同步参考床垫离床状态、毫米波雷达静止时间、可穿戴设备活动变化和历史健康档案，从而判断是否进入高风险或紧急事件队列。

\subsection{面向健康与心理风险的个体化行为画像}

系统不只关注跌倒、烟火等显性安全事件，还关注久卧未动、夜间离床频繁、活动量下降、情绪低落、孤独风险、焦虑风险和睡眠异常等健康与心理风险。项目通过老人历史作息、活动区域、护理记录和近期异常次数构建个体化行为画像，使系统能够识别“相对于本人历史状态的异常”，而不是只依赖统一阈值判断。

\subsection{视觉识别与姿态估计结合}

项目结合 YOLO11n 和 YOLO11n-pose，将目标检测与姿态估计共同用于疗养中心场景。目标检测用于识别人员、床位、轮椅、区域目标和环境异常，姿态估计用于判断疑似跌倒、异常姿态、长时间静止和冲突动作。二者结合后，可以为风险融合层提供更高置信度的感知输入。

\subsection{可解释风险研判}

系统不仅输出风险等级，还输出触发原因、风险评分、置信度、知识库依据和建议动作。例如，系统可以说明某事件被判定为高风险的原因是“夜间离床超过阈值、门磁未触发回房记录、床垫传感器持续空床、老人近三日睡眠质量下降”。这种可解释输出便于护理员确认现场情况，也便于医生理解长期健康趋势。

\subsection{从预警到处置的护理业务闭环}

项目覆盖风险发现、智能研判、派单调度、应急流程、事件记录、统计看板和知识库辅助，实现疗养中心后台管理闭环。系统不是孤立告警工具，而是将 AI 预警嵌入护理工作流，使事件能够被接收、响应、处理、归档和复盘。

\subsection{护理知识库辅助决策}

针对跌倒、久卧、夜间离床、心理风险、无人看护等事件，系统可生成处置步骤、注意事项和现场护理建议。知识库既可作为应急流程依据，也可作为护理培训和事件复盘材料，提升机构处置规范性和经验沉淀能力。

\section{项目的应用前景和社会价值}

\subsection{应用前景}

安养智巡适用于老年人疗养中心、护理院、医养结合机构、康复照护中心和长期照护病区等场景。系统可以部署为机构后台管理平台，与摄像头、床垫传感器、门磁、毫米波雷达、可穿戴设备和现有护理记录系统对接，形成统一风险管理入口。对于中小型机构，系统可优先部署核心页面和事件闭环模块；对于大型连锁机构，系统可扩展为多院区、多角色、多权限和多看板的运营管理平台。

\subsection{社会价值}

第一，系统能够提高老年人风险发现的及时性。通过多模态连续感知，系统可以在护理员巡房间隔之外持续捕捉异常信号，减少跌倒未响应、夜间离床未归和久卧未动等事件的延迟发现。

第二，系统能够提升护理资源调度效率。风险调度模块根据风险等级、等待时长、护理员状态和区域位置排序，使有限护理资源优先投向更紧急、更需要现场处理的事件。

第三，系统能够促进健康和心理风险的早期干预。行为画像和健康监护模块使机构能够关注活动量下降、睡眠异常、孤独风险和焦虑风险等长期变化，为医生和护理主管提供更完整的观察依据。

第四，系统能够推动养老机构数字化管理。事件管理、数据看板和知识库工作台将分散的护理记录转化为可统计、可追踪、可复盘的数据资产，有助于改进护理流程、设备配置和人员排班。

\section{团队分工与协同实现}

本项目由 3 名成员协同完成。团队在系统方案、算法服务、后台页面、接口联调、测试数据和部署验证等环节保持共同讨论与交叉验证，分工围绕真实康养机构后台业务展开，重点保证系统能够形成从风险发现到事件归档的完整链路。

队员 A 主要负责系统总体架构设计、多模态融合预警流程设计、核心业务闭环设计、报告撰写和系统功能整合。在项目推进过程中，队员 A 同时参与风险评分逻辑、页面业务含义梳理和系统部署方案设计，使系统从技术模块组合进一步形成可解释、可追踪的业务闭环。

队员 B 主要负责后台管理系统开发，包括首页、实时监测、健康监护、风险调度、行为画像、应急流程、事件管理、数据看板和知识库工作台等页面。队员 B 在实现过程中重点处理多角色入口、页面状态、业务数据联动和事件流转体验，使系统能够符合护理员、值班主管、机构管理者和医生的使用习惯。

队员 C 主要负责视觉识别服务、YOLO 模型接入、风险规则设计、后端接口联调、测试数据构建和系统部署验证。队员 C 同时参与视觉识别结果与传感器状态融合、事件队列生成和运行测试，保证算法输出能够稳定进入后端服务和业务处置流程。

总体来看，三名成员均参与了创新设计和核心实现。队员 A 对系统架构、融合预警方案和业务闭环完整性贡献突出；队员 B 对后台管理系统的人机交互、功能页面和业务流程落地贡献突出；队员 C 对视觉识别、规则研判、后端接口和部署验证贡献突出。团队协同方式保证了项目不仅具备页面呈现能力，也具备服务接口、算法接入、规则研判、事件闭环和部署验证能力。

\section{项目存在的问题以及今后的改进方向}

\subsection{存在的问题}

第一，真实场景数据规模仍需进一步扩大。当前系统已完成原型部署与模拟运行验证，主要业务流程已经能够连贯运行，但不同疗养中心的老人状态、空间布局、护理制度和设备布设差异较大，后续仍需更多真实数据持续优化模型。

第二，不同机构的泛化能力仍需验证。疗养中心的房间结构、走廊宽度、摄像头角度、床位布局、遮挡情况和传感器型号存在差异，视觉模型和规则阈值需要根据场景进行适配。

第三，多模态数据接入标准仍需进一步完善。不同厂商设备的数据格式、通信协议和采样频率不一致，后续需要建立更规范的数据接入层和设备抽象接口。

第四，健康和心理风险识别需要长期数据支持。心理风险、孤独风险和慢性健康趋势并非短时间即可准确判断，需要结合更长周期的行为画像、护理记录和医生反馈持续校准。

第五，隐私保护、边缘部署和权限管理仍有加强空间。摄像头和健康数据具有敏感性，后续需要进一步完善数据脱敏、边缘推理、访问控制、日志审计和合规管理设计。

\subsection{改进方向}

后续将从七个方向持续优化：第一，接入更多真实摄像头视频流和传感器数据，扩大跌倒、久卧、夜间离床和心理风险样本；第二，优化视觉识别、姿态估计和行为偏离度模型，提高复杂遮挡、弱光和多人场景下的稳定性；第三，引入边缘计算设备，在机构本地完成关键视觉推理，降低延迟并减少敏感数据外传；第四，完善数据库、权限系统、日志审计和机构级多角色管理；第五，引入大模型护理助手，提高知识库问答、处置建议生成和事件复盘摘要质量；第六，建立长期健康画像和心理风险趋势预测模型；第七，加强数据脱敏、访问控制和隐私合规设计，使系统更适合真实养老机构长期运行。

\phantomsection
\addcontentsline{toc}{section}{参考文献}
\begin{thebibliography}{99}
\songti\zihao{5}
\setlength{\itemsep}{0pt}
\setlength{\parsep}{0pt}
\setlength{\parskip}{0pt}

\bibitem{zhang2020smart}
Zhang Q, Li M, Wu Y. Smart home for elderly care: development and challenges in China[J]. \textit{BMC Geriatrics}, 2020, 20:318.

\bibitem{hung2023smart}
Hung J. Smart Elderly Care Services in China: Challenges, Progress, and Policy Development[J]. \textit{Sustainability}, 2023, 15(1):178.

\bibitem{stavropoulos2020iot}
Stavropoulos T G, Papastergiou A, Mpaltadoros L, Nikolopoulos S, Kompatsiaris I. IoT Wearable Sensors and Devices in Elderly Care: A Literature Review[J]. \textit{Sensors}, 2020, 20(10):2826.

\bibitem{dong2017wearable}
Dong L, et al. A Review of Wearable Technologies for Elderly Care that Can Accurately Track Indoor Position, Recognize Physical Activities and Monitor Vital Signs in Real Time[J]. \textit{Sensors}, 2017, 17(2):341.

\bibitem{mubashir2013survey}
Mubashir M, Shao L, Seed L. A survey on fall detection: Principles and approaches[J]. \textit{Neurocomputing}, 2013, 100:144--152.

\bibitem{igual2013challenges}
Igual R, Medrano C, Plaza I. Challenges, issues and trends in fall detection systems[J]. \textit{BioMedical Engineering OnLine}, 2013, 12:66.

\bibitem{redmon2016yolo}
Redmon J, Divvala S, Girshick R, Farhadi A. You Only Look Once: Unified, Real-Time Object Detection[C]//\textit{Proceedings of the IEEE Conference on Computer Vision and Pattern Recognition}. 2016:779--788.

\bibitem{cao2017realtime}
Cao Z, Simon T, Wei S E, Sheikh Y. Realtime Multi-Person 2D Pose Estimation Using Part Affinity Fields[C]//\textit{Proceedings of the IEEE Conference on Computer Vision and Pattern Recognition}. 2017:7291--7299.

\bibitem{ultralytics2024yolo11}
Ultralytics. Ultralytics YOLO11 Documentation[EB/OL]. 2024.

\bibitem{baltrusaitis2019multimodal}
Baltrušaitis T, Ahuja C, Morency L P. Multimodal Machine Learning: A Survey and Taxonomy[J]. \textit{IEEE Transactions on Pattern Analysis and Machine Intelligence}, 2019, 41(2):423--443.

\bibitem{liang2022foundations}
Liang P P, Zadeh A, Morency L P. Foundations and Trends in Multimodal Machine Learning: Principles, Challenges, and Open Questions[Z]. arXiv:2209.03430, 2022.

\bibitem{holtlunstad2015loneliness}
Holt-Lunstad J, Smith T B, Baker M, Harris T, Stephenson D. Loneliness and Social Isolation as Risk Factors for Mortality: A Meta-Analytic Review[J]. \textit{Perspectives on Psychological Science}, 2015, 10(2):227--237.

\bibitem{sequeira2021digital}
Sequeira L, Perrotta S, LaGrassa J, et al. Digital health tools for the passive monitoring of depression: a systematic review of methods[J]. \textit{npj Digital Medicine}, 2021, 4:143.

\bibitem{kawamoto2005cdss}
Kawamoto K, Houlihan C A, Balas E A, Lobach D F. Improving clinical practice using clinical decision support systems: a systematic review of trials to identify features critical to success[J]. \textit{BMJ}, 2005, 330:765.

\bibitem{bright2012cdss}
Bright T J, Wong A, Dhurjati R, et al. Effect of Clinical Decision-Support Systems: A Systematic Review[J]. \textit{Annals of Internal Medicine}, 2012, 157(1):29--43.

\bibitem{gunning2019xai}
Gunning D, Aha D. DARPA's Explainable Artificial Intelligence Program[J]. \textit{AI Magazine}, 2019, 40(2):44--58.

\bibitem{lundberg2017shap}
Lundberg S M, Lee S I. A Unified Approach to Interpreting Model Predictions[C]//\textit{Advances in Neural Information Processing Systems}. 2017, 30.

\end{thebibliography}

\end{document}
